\chapter{Herleitungen}

\section{Einleitung}

Dieser Anhang enthält die vollständigen Herleitungen der zentralen Gleichungen der IWT und der WDBT+.

Die Herleitungen folgen aus den Axiomen der IWT und den Prinzipien der diskreten Variationsrechnung.

\section{Herleitung von Q}

Aus den Prinzipien der Informationserhaltung und minimaler Spannung folgt \( Q \).

\subsection{Der Ansatz}

Wir betrachten ein Funktional, das die Spannung im Netzwerk misst:

\[
\mathcal{S} = \sum_k |\nabla I_k|^2
\]

Die Informationserhaltung wird als Nebenbedingung eingeführt:

\[
\sum_k |I_k|^2 = \text{const}
\]

\subsection{Die Lagrange-Multiplikator-Methode}

Wir minimieren die Spannung unter der Nebenbedingung:

\[
\mathcal{L} = \sum_k |\nabla I_k|^2 + \lambda \left( \sum_k |I_k|^2 - \text{const} \right)
\]

Die Variation nach \( I_k \) ergibt:

\[
\delta \mathcal{L} = 0
\]

Daraus folgt die Euler-Lagrange-Gleichung:

\[
\Delta^2 I_k = \lambda \cdot I_k
\]

\subsection{Die Lösung}

Die Lösung hat die Form:

\[
I_k = \sqrt{\rho_k} \cdot e^{i \phi_k}
\]

Einsetzen in die Euler-Lagrange-Gleichung und Trennung von Real- und Imaginärteil ergibt:

\[
\frac{\partial \rho}{\partial t} + \nabla \cdot (\rho \vec{v}) = 0
\]

und:

\[
\frac{\partial S}{\partial t} + \frac{|\nabla S|^2}{2m} + V + Q = 0
\]

mit:

\[
Q = -\frac{\hbar^2}{2m} \frac{\nabla^2 \sqrt{\rho}}{\sqrt{\rho}}
\]

\subsection{Die Identifikation}

\( Q \) ist das Bohm-Potential der De-Broglie-Bohm-Theorie.

Es ist die Quelle der Selbstorganisation.

Es ist die bindende Kraft.

\section{Herleitung der DBT aus Q}

Aus \( Q \) folgt die Führungsgleichung der DBT:

\[
\vec{v} = \frac{\hbar}{m} \nabla \phi
\]

Die Schrödinger-Gleichung folgt aus der Hamilton-Jacobi-Gleichung mit dem Quantenpotential \( Q \):

\[
i\hbar \frac{\partial \psi}{\partial t} = -\frac{\hbar^2}{2m} \nabla^2 \psi + V \psi
\]

wobei \( \psi = \sqrt{\rho} e^{i\phi} \).

\section{Herleitung der WG aus Q}

Aus \( Q \) folgt die Weber-Gravitation.

\subsection{Die Kraft}

Die Weber-Gravitation lautet:

\[
\vec{F}_{\text{WG}} = -\frac{GMm}{r^2} \left[ 1 - \frac{\dot{r}^2}{c^2} + \frac{1}{2} \frac{r\ddot{r}}{c^2} \right] \hat{r}
\]

Aus \( Q \) folgt:

\[
\vec{F}_{\text{WG}} = -\nabla Q
\]

mit:

\[
Q = -\frac{GM}{r}
\]

\subsection{Die Periheldrehung}

Die Bahngleichung der WG in Binet-Form:

\[
u'' + u = \frac{GM}{h^2} + \frac{3GM}{c^2} u^2
\]

Die Lösung ergibt:

\[
\Delta \phi = \frac{6\pi GM}{c^2 a (1 - e^2)}
\]

Für Merkur ergibt sich \( \Delta \phi \approx 42,98'' \) pro Jahrhundert.

\section{Herleitung der WED aus Q}

Aus \( Q \) folgt die Weber-Elektrodynamik.

\subsection{Die Kraft}

Die Weber-Elektrodynamik lautet:

\[
\vec{F}_{\text{WED}} = \frac{q_1 q_2}{4\pi \varepsilon_0 r^2} \left[ 1 - \frac{\dot{r}^2}{c^2} + 2 \frac{r\ddot{r}}{c^2} \right] \hat{r}
\]

Aus \( Q \) folgt:

\[
\vec{F}_{\text{WED}} = -\nabla Q
\]

mit:

\[
Q = \frac{q}{4\pi \varepsilon_0 r}
\]

\section{Herleitung der Unschärferelation}

Aus der diskreten Zeit \( T > 0 \) folgt die Unschärferelation.

\subsection{Bandbegrenzung}

Die diskrete Zeit erzwingt ein bandbegrenztes Frequenzspektrum:

\[
\omega_{\text{max}} = \frac{\pi}{T}
\]

\subsection{Bandbreiten-Unschärfe}

Für jedes bandbegrenzte Signal gilt:

\[
\Delta t \cdot \Delta \omega \geq \frac{1}{2}
\]

\subsection{Energie-Zeit-Unschärfe}

Mit \( E = \hbar \omega \) folgt:

\[
\Delta t \cdot \Delta E \geq \frac{\hbar}{2}
\]

\subsection{Orts-Impuls-Unschärfe}

Mit \( p = \hbar k \) folgt:

\[
\Delta x \cdot \Delta p \geq \frac{\hbar}{2}
\]

Die Unschärferelation ist eine zwingende Konsequenz der diskreten Zeit.

\section{Herleitung der fraktalen Dimension}

Die fraktale Dimension folgt aus der selbstähnlichen Packung von Dodekaedern.

\subsection{Der Skalierungsfaktor}

Aus der Dodekaeder-Geometrie folgt:

\[
s = 2 + \phi \approx 3,618
\]

\subsection{Der Vermehrungsfaktor}

Die Anzahl der Ecken eines Dodekaeders ist \( V = 20 \). Mit dem Goldenen Schnitt:

\[
N = 20\phi \approx 32,36
\]

\subsection{Die fraktale Dimension}

\[
D = \frac{\ln N}{\ln s} = \frac{\ln(20\phi)}{\ln(2 + \phi)} \approx 2,704
\]

\section{Zusammenfassung}

\begin{itemize}
\item \( Q \) folgt aus Informationserhaltung und Variationsprinzip.
\item DBT, WG und WED folgen aus \( Q \).
\item Die Unschärferelation folgt aus der diskreten Zeit.
\item Die fraktale Dimension folgt aus der Dodekaeder-Geometrie.
\end{itemize}